% Tóm tắt (khoảng 200 từ). Mọi số đo được dùng macro của generated/numbers.tex.
\begin{infobox}{Tóm tắt}
Bài LTSF-Linear cho thấy các mô hình tuyến tính DLinear và NLinear dự báo chuỗi thời gian dài tốt hơn nhiều mô hình Transformer, nhưng mọi so sánh đều qua huấn luyện bằng SGD. Vì vậy chưa rõ ưu thế đến từ kiến trúc hay từ cách huấn luyện. Chúng tôi thay SGD bằng nghiệm tối ưu tất định: nghiệm dạng đóng cho MSE (OLS, ridge) và bình phương tối thiểu có trọng số lặp lại (IRLS) cho MAE và Huber. Ma trận thí nghiệm gồm 3 mô hình, 3 hàm mục tiêu, 3 dataset ETT và 4 horizon; $\lambda$ chỉ chọn trên tập kiểm định. Có ba kết quả chính.
(1)~Tại nghiệm tối ưu, DLinear trùng Linear khi không phạt, và chỉ tách khỏi Linear khi có weight decay. Trên mọi dataset và horizon, MSE test của DLinear và Linear tại $\lambda^*$ chênh nhau không quá \nmBThreeDLinMinusLinAbsMax. NLinear tách khỏi Linear ở mọi $\lambda$, vì ràng buộc tổng hàng bằng $1$; ràng buộc này có lợi rõ trên ETTh2.
(2)~Huấn luyện bằng MAE cho MSE test thấp hơn huấn luyện bằng MSE ở \nmBThreeMaeBeatsMseTotalCount/\nmBThreeCellCount{} ô, và thấp nhất trong ba mục tiêu ở \nmBThreeBestMAEMseCount{} ô. Huber nằm giữa MSE và MAE ở \nmBThreeHuberRankTwoMseCount/\nmBThreeCellCount{} ô. MAE bền vững khi outlier chỉ nằm ở đích, nhưng mất lợi thế khi outlier vừa nhiều vừa lớn và nằm cả trong đầu vào.
(3)~Sau khi tinh chỉnh learning rate, SGD vẫn kém nghiệm dạng đóng trên test ở \nmTunedWorseSigCount/\nmBThreeCellCount{} ô (có ý nghĩa thống kê). Ngoại lệ chỉ có trên ETTh2. Khi cả hai mô hình được tinh chỉnh learning rate, DLinear không còn tốt hơn Linear có ý nghĩa ở ô nào.
\end{infobox}
